% Relatorio de sistema: inventario do HiGTree/HiGFlow.
% Mesma convencao dos outros relatorios desta pasta: os NUMEROS vivem em
% dados/*.dat e sao lidos por pgfplotstable, nao digitados no texto.
\documentclass[11pt,a4paper]{article}
\usepackage[T1]{fontenc}
\usepackage[utf8]{inputenc}
\usepackage[brazil]{babel}
\usepackage{lmodern}
\usepackage{microtype}
\usepackage{amsmath,amssymb}
\usepackage{booktabs}
\usepackage{siunitx}
\usepackage{pgfplotstable}
\usepackage[hidelinks]{hyperref}
\usepackage[margin=2.7cm]{geometry}
\pgfplotsset{compat=1.18}
\sisetup{output-decimal-marker={,}, group-separator={.}}
\newcommand{\dat}[1]{dados/#1}

\title{O sistema HiGTree/HiGFlow:\\ invent\'ario dos m\'odulos existentes}
\author{Relat\'orio t\'ecnico}
\date{Setembro de 2026}

\begin{document}
\maketitle

\begin{abstract}
Invent\'ario do que existe no c\'odigo do HiGTree/HiGFlow: 75 arquivos-fonte e
cerca de 76 mil linhas, distribu\'idas entre a biblioteca de malha, o solver, os
adaptadores compartilhados e dezessete exemplos execut\'aveis. Os n\'umeros vem
de contagem direta da \'arvore.

O levantamento torna vis\'iveis tr\^es coisas que a leitura do c\'odigo por
partes n\~ao mostra: a entrada e sa\'ida \'e o maior m\'odulo isolado do sistema,
com \SI{20}{\percent} do HiGFlow; os modelos reol\'ogicos ocupam tr\^es vezes
mais linhas que a infraestrutura do solver; e a combinat\'oria de f\'isicas
est\'a mais \emph{declarada} que realizada --- dois dos cinco acoplamentos t\^em
menos de trinta linhas e apenas despacham.

Este \'e um invent\'ario, e n\~ao uma avalia\c{c}\~ao: a presen\c{c}a de um
m\'odulo n\~ao diz que ele esteja correto, verificado ou em uso.
\end{abstract}

\tableofcontents

\section{Introdu\c{c}\~ao}

\subsection{O que \'e o sistema}

O HiGFlow \'e um simulador de escoamentos incompress\'iveis constru\'ido sobre o
HiGTree, uma biblioteca de malhas em \'arvore \textbf{n\~ao graduada} com
interpola\c{c}\~ao sem malha. A funda\c{c}\~ao do m\'etodo est\'a em Sousa et
al.~\cite{sousa2019}: diferen\c{c}as finitas em \'arvore onde c\'elulas vizinhas
podem diferir de mais de um n\'ivel de refino, com a interpola\c{c}\~ao por
m\'inimos quadrados ponderados cobrindo o que o estêncil regular n\~ao alcan\c{c}a.

A escolha de n\~ao exigir gradua\c{c}\~ao \'e o que separa o sistema da maioria
dos c\'odigos de AMR, e ela aparece em toda parte: no
\texttt{domain.c}, que monta estênceis sem supor vizinhan\c{c}a uniforme; no
\texttt{wls.c}, que \'e a interpola\c{c}\~ao que torna isso poss\'ivel; e no custo
--- cada avalia\c{c}\~ao fora do estêncil regular resolve um sistema local.

\subsection{A linhagem}

O sistema n\~ao come\c{c}ou em 2019. Ele \'e a terceira gera\c{c}\~ao de uma
linha de simuladores do mesmo grupo:

\begin{itemize}
  \item \textbf{GENSMAC}~\cite{tome1994} --- marcadores e c\'elulas para
    superf\'icie livre em dom\'inios gerais, 1994;
  \item \textbf{Freeflow}~\cite{castelo2000} --- o sistema integrado 3D que
    embrulhou o GENSMAC em ambiente de simula\c{c}\~ao, 2000, com a vers\~ao
    axissim\'etrica em~\cite{tome2000};
  \item \textbf{o front-tracking do pr\'oprio grupo}~\cite{sousa2004} --- a
    combina\c{c}\~ao de rastreamento e captura de frente para multifluido 3D,
    2004, que \'e o antecedente direto do m\'odulo de front-tracking descrito
    adiante;
  \item \textbf{HiGTree/HiGFlow}~\cite{sousa2019} --- a malha em \'arvore n\~ao
    graduada, 2019.
\end{itemize}

Ler a lista na ordem explica escolhas do c\'odigo que de outro modo pareceriam
arbitr\'arias --- a presen\c{c}a simult\^anea de VOF e front-tracking, por
exemplo, reproduz a coexist\^encia que~\cite{sousa2004} j\'a propunha.

\subsection{O que este relat\'orio \'e, e o que n\~ao \'e}

\'E um \textbf{invent\'ario}: o que existe no c\'odigo, em quantas linhas, e qual
o papel de cada m\'odulo. Os n\'umeros vem de contagem direta da \'arvore, n\~ao
de estimativa.

N\~ao \'e uma avalia\c{c}\~ao de corre\c{c}\~ao. Dizer que um m\'odulo existe
n\~ao \'e dizer que ele est\'a certo, nem que est\'a verificado, nem que \'e
usado. Onde houver verifica\c{c}\~ao conhecida, ela est\'a marcada; onde n\~ao
houver, a aus\^encia tamb\'em.

\subsection{Sobre as refer\^encias, e o limite desta busca}

As entradas da bibliografia foram conferidas contra o registro do editor via API
do Crossref, e n\~ao contra mem\'oria.

\paragraph{A p\'os-gradua\c{c}\~ao sobre o sistema.} A busca localizou
\textbf{uma} disserta\c{c}\~ao da USP com o sistema no t\'itulo ---
Godoi~\cite{godoi2022}, ICMC/S\~ao Carlos, aprovada em novembro de 2022, sobre
modelos viscoel\'asticos com viscosidade vari\'avel. Ela corresponde diretamente
a um m\'odulo desta \'arvore, o
\texttt{hig-flow-step-viscoelastic-variable-viscosity.c}, com suas 2\,435 linhas.

\textbf{E aqui o limite precisa ser dito, porque ele \'e grande.} O Crossref
indexa metadados --- t\'itulo, autores, resumo ---, e n\~ao texto integral. Uma
tese que \emph{use} o HiGFlow sem nome\'a-lo no t\'itulo n\~ao aparece nessa
busca. Dado que o sistema tem m\'odulos para tixotropia, bandas de cisalhamento,
suspens\~oes, eletro-osmose e viscoelasticidade integral --- cada um com milhares
de linhas, e cada um o tipo de assunto que rende uma disserta\c{c}\~ao ---, \'e
quase certo que existam outras. Elas n\~ao est\~ao aqui porque \textbf{n\~ao foram
encontradas}, e n\~ao porque n\~ao existam.

Listar refer\^encias prov\'aveis sem conferi-las seria pior que omiti-las: um
relat\'orio de sistema \'e consultado justamente para saber o que citar. O
caminho para completar esta se\c{c}\~ao \'e uma busca de texto integral no
reposit\'orio de teses da USP, que est\'a fora do alcance das ferramentas usadas
aqui.

\paragraph{Os contribuidores do reposit\'orio.} O hist\'orico Git cobre de
setembro de 2020 a setembro de 2026, com 358 \emph{commits}. Os nomes que
aparecem s\~ao Antonio Castelo Filho, Juniormar Organista, Daniel Garcia,
Johnatas e Pedro Vin\'icius. Isso \'e autoria de \emph{commit}, que n\~ao
coincide necessariamente com autoria dos m\'odulos nem com as teses ---
registro o dado, n\~ao a infer\^encia.

\section{HiGTree: a camada de malha}

Trinta e tr\^es fontes, 15\,869 linhas. \'E a biblioteca sobre a qual tudo o mais
se ap\'oia, e ela n\~ao sabe nada de Navier--Stokes.

\begin{adjustbox}{max width=\linewidth}
\pgfplotstabletypeset[
  string type,
  columns/modulo/.style={column name={m\'odulo}},
  columns/linhas/.style={column name={linhas}},
  columns/papel/.style={column name={papel}},
  every head row/.style={before row=\toprule, after row=\midrule},
  every last row/.style={after row=\bottomrule},
]{\dat{higtree-modulos.dat}}
\end{adjustbox}

\paragraph{A divis\~ao que mais importa \'e \texttt{domain} contra
\texttt{pdomain}.} O primeiro \'e o dom\'inio serial --- c\'elulas, facetas,
estênceis, mapeadores. O segundo \'e o distribu\'ido, e \'e onde vivem a
\emph{franja} (as c\'opias de c\'elulas de outro \emph{rank}), a
sincroniza\c{c}\~ao e os tipos derivados de MPI. Juntos s\~ao
5\,115 linhas, quase um ter\c{c}o da biblioteca: a distribui\c{c}\~ao \'e o
assunto mais caro dela.

\paragraph{Solvers lineares: cinco implementa\c{c}\~oes atr\'as de uma
interface.} PETSc, HYPRE, SOR, centralizado e um de depura\c{c}\~ao que
simplesmente grava o sistema. O de PETSc \'e o usado na pr\'atica. Existe
tamb\'em um \texttt{solver-\allowbreak{}petsc.\allowbreak{}camila.\allowbreak{}c}, variante paralela do principal ---
duas c\'opias de 576 e 555 linhas que divergir\~ao na primeira mudan\c{c}a de
protocolo.

\paragraph{\texttt{wls.c} \'e o cora\c{c}\~ao do m\'etodo.} Sem os m\'inimos
quadrados ponderados, a \'arvore n\~ao graduada n\~ao fecha: \'e ele que produz
valor onde o estêncil regular n\~ao existe. As 624 linhas s\~ao pequenas para o
peso que carregam.

\section{HiGFlow: o solver}

Quarenta e dois fontes, 54\,457 linhas --- tr\^es vezes e meia a HiGTree. A maior
parte disso s\~ao \emph{modelos}, n\~ao infraestrutura.

\subsection{Um m\'odulo de passo por tipo de escoamento}

O \texttt{flowtype} da configura\c{c}\~ao seleciona qual passo roda:

\pgfplotstabletypeset[
  string type,
  columns/tipodeescoamento/.style={column name={\texttt{flowtype}}},
  columns/modulo/.style={column name={m\'odulo}},
  columns/linhas/.style={column name={linhas}},
  every head row/.style={before row=\toprule, after row=\midrule},
  every last row/.style={after row=\bottomrule},
]{\dat{higflow-passos.dat}}

S\~ao 19\,900 linhas s\'o em passos de escoamento. O maior deles ---
suspens\~oes, com 3\,413 --- \'e maior que \texttt{domain.c}, o n\'ucleo da
HiGTree.

\subsection{E os acoplamentos entre eles}

\pgfplotstabletypeset[
  string type,
  columns/modulo/.style={column name={m\'odulo}},
  columns/linhas/.style={column name={linhas}},
  columns/acopla/.style={column name={acopla}},
  every head row/.style={before row=\toprule, after row=\midrule},
  every last row/.style={after row=\bottomrule},
]{\dat{higflow-acoplados.dat}}

Os dois \'ultimos t\^em 31 e 28 linhas: s\~ao despachantes, n\~ao
implementa\c{c}\~oes. A combinat\'oria de f\'isicas est\'a declarada mais do que
realizada, e vale saber disso antes de contar com ela.

\subsection{Modelos constitutivos}

\pgfplotstabletypeset[
  string type,
  columns/familia/.style={column name={fam\'ilia}},
  columns/modelos/.style={column name={modelos}},
  every head row/.style={before row=\toprule, after row=\midrule},
  every last row/.style={after row=\bottomrule},
]{\dat{modelos.dat}}

\subsection{Discretiza\c{c}\~oes, por eixo}

\pgfplotstabletypeset[
  string type,
  columns/eixo/.style={column name={eixo}},
  columns/opcoes/.style={column name={op\c{c}\~oes}},
  every head row/.style={before row=\toprule, after row=\midrule},
  every last row/.style={after row=\bottomrule},
]{\dat{discretizacoes.dat}}

O produto cartesiano dessas op\c{c}\~oes \'e grande, e nada garante que todas as
combina\c{c}\~oes tenham sido exercitadas --- esta \'e uma lista do que o c\'odigo
\emph{aceita}, n\~ao do que foi verificado.

\subsection{Tratamento de interface}

\pgfplotstabletypeset[
  string type,
  columns/familia/.style={column name={fam\'ilia}},
  columns/modulos/.style={column name={m\'odulos}},
  columns/papel/.style={column name={papel}},
  every head row/.style={before row=\toprule, after row=\midrule},
  every last row/.style={after row=\bottomrule},
]{\dat{higflow-interface.dat}}

Onze m\'odulos de VOF, contra um de front-tracking e um de fronteira imersa. A
assimetria conta a hist\'oria do sistema: o VOF \'e o caminho maduro, com
variantes de reconstru\c{c}\~ao e de curvatura acumuladas ao longo de anos; o
front-tracking \'e recente e tem uma implementa\c{c}\~ao s\'o.

\subsection{Entrada e sa\'ida: o m\'odulo mais improv\'avel}

\texttt{hig-flow-io.c} tem \textbf{10\,972 linhas} --- \SI{20}{\percent} de todo o
HiGFlow, e mais que a soma de \texttt{domain.c} e \texttt{pdomain.c}. Ele l\^e as
configura\c{c}\~oes em YAML, grava e rel\^e o estado, e escreve VTK. Cada
combina\c{c}\~ao de \texttt{flowtype} com dimens\~ao tem o seu par de
leitura/escrita, e \'e da\'i que vem o tamanho.

\'E o m\'odulo que ningu\'em planeja e que cresce sozinho.

\section{Exemplos e c\'odigo compartilhado}

\subsection{Os dezessete casos}

\pgfplotstabletypeset[
  string type,
  columns/exemplo/.style={column name={exemplo}},
  columns/o-que-exercita/.style={column name={o que exercita}},
  every head row/.style={before row=\toprule, after row=\midrule},
  every last row/.style={after row=\bottomrule},
]{\dat{exemplos.dat}}

Cada exemplo \'e um programa pr\'oprio, com o seu \texttt{main} e o seu
\texttt{Makefile}, que linka os objetos de \texttt{higflow/src} de que precisa.
N\~ao h\'a biblioteca \'unica do HiGFlow: h\'a uma cole\c{c}\~ao de objetos e
dezessete formas de combin\'a-los.

\subsection{\texttt{examples-common}: o que os exemplos compartilham}

Cinco mil seiscentas e cinquenta e oito linhas fora de \texttt{src}, e \'e onde
vive o c\'odigo que \'e demais para um exemplo e de menos para a biblioteca:

\begin{itemize}
  \item \texttt{utilities-2d.c} (2\,024 linhas) --- condi\c{c}\~oes de contorno e
    iniciais compartilhadas;
  \item \texttt{front-tracking.c} (995) --- o adaptador entre o n\'ucleo
    lagrangeano e o solver;
  \item \texttt{malha-adaptativa.c} (789) --- crit\'erio de refino e
    reconstru\c{c}\~ao, compartilhado entre VOF e front-tracking;
  \item \texttt{malha-t8.cxx} (442) --- as fontes de malha do t8code, compiladas
    apenas quando \texttt{T8CODE} \'e passado ao \texttt{make};
  \item \texttt{vof-geometry-2d.c} e \texttt{eo-droplet-2d.c} --- geometrias
    iniciais espec\'ificas.
\end{itemize}

\subsection{Integra\c{c}\~ao com o t8code}

O t8code entra como \emph{fonte de malha} opcional, com dupla ades\~ao: em tempo
de compila\c{c}\~ao (\texttt{T8CODE=<prefixo>}) e em tempo de execu\c{c}\~ao
(\texttt{HIGFLOW\_MALHA=}). Sem as duas, o caminho \'e o mtree, lendo o
\texttt{.amr}.

Duas limita\c{c}\~oes, ambas verificadas no c\'odigo:

\begin{enumerate}
  \item o caminho \textbf{multif\'asico} n\~ao consulta a fonte de malha ---
    \texttt{higflow\_partition\_domain\_multiphase} l\^e o \texttt{.amr} sempre.
    A fronteira imersa funciona com t8code porque \'e monof\'asica e passa por
    outra fun\c{c}\~ao;
  \item as fontes do t8code \textbf{recusam} \texttt{.amr} multin\'ivel, com
    mensagem expl\'icita (``recusar \'e melhor que aproximar''). Como o
    crit\'erio adaptativo escreve exatamente isso, integrar o t8code ao caso
    adaptativo exige uma fonte que \emph{leia} a lista de refino --- que o
    \texttt{higio\_amr\_info} carrega e entrega, e que hoje \'e descartada em
    favor s\'o da caixa envolvente.
\end{enumerate}

\section{O que o invent\'ario mostra}

\subsection{Onde est\~ao as linhas}

\pgfplotstabletypeset[
  string type,
  columns/camada/.style={column name={camada}},
  columns/arquivos/.style={column name={arquivos}},
  columns/linhas/.style={column name={linhas}},
  columns/papel/.style={column name={papel}},
  every head row/.style={before row=\toprule, after row=\midrule},
  every last row/.style={after row=\bottomrule},
]{\dat{tamanhos.dat}}

Tr\^es observa\c{c}\~oes que a contagem torna dif\'iceis de ignorar.

\paragraph{A E/S \'e o maior m\'odulo isolado do sistema.} \texttt{hig-flow-io.c}
sozinho tem \SI{20}{\percent} do HiGFlow e mais linhas que o dom\'inio serial e o
distribu\'ido somados. A causa \'e estrutural: cada \texttt{flowtype} $\times$
dimens\~ao $\times$ opera\c{c}\~ao tem o seu par de leitura e escrita, e o
crescimento \'e multiplicativo.

\paragraph{Os modelos dominam o solver.} Dezenove mil e novecentas linhas em
passos de escoamento, contra cerca de seis mil de infraestrutura (\texttt{kernel},
\texttt{discret}, \texttt{bc}, \texttt{ic}, \texttt{eval}, \texttt{terms}). O
sistema \'e, em volume, uma cole\c{c}\~ao de modelos reol\'ogicos com um solver
de Navier--Stokes dentro.

\paragraph{A combinat\'oria est\'a declarada mais do que realizada.} H\'a nove
\texttt{flowtype}, cinco acoplamentos, doze modelos constitutivos, seis
discretiza\c{c}\~oes temporais e duas proje\c{c}\~oes. O produto \'e enorme, e
dois dos acoplamentos t\^em menos de trinta linhas. Uma lista de capacidades
lida como cat\'alogo superestima o que est\'a exercitado.

\subsection{O que este invent\'ario n\~ao diz}

\begin{itemize}
  \item \textbf{nada sobre corre\c{c}\~ao}. Contagem de linhas n\~ao \'e
    verifica\c{c}\~ao, e a presen\c{c}a de um m\'odulo n\~ao diz se ele foi
    validado contra or\'aculo algum;
  \item \textbf{nada sobre uso}. N\~ao foi medido quais m\'odulos os exemplos
    de fato exercitam, e h\'a sinais de duplica\c{c}\~ao --- duas variantes do
    solver de PETSc, dois de \texttt{solver.c};
  \item \textbf{nada sobre paralelismo}. Quais m\'odulos funcionam em $n_p>1$
    \'e outra pergunta, e a resposta varia por m\'odulo;
  \item \textbf{a p\'os-gradua\c{c}\~ao associada est\'a incompleta}, pelas
    raz\~oes da introdu\c{c}\~ao: a busca alcan\c{c}a metadados, n\~ao texto
    integral.
\end{itemize}


\bibliographystyle{plain}
\bibliography{referencias}

\end{document}
